

\paragraph{Goal.}
We ask a mechanistic question: when a vision--language model $L$ is adapted
to a shifted image distribution via an adapter $A$ to produce $\hat{L}$, does
the feature basis learned by a sparse autoencoder $S_L$ on the base model
remain a faithful decomposition of $\hat{L}$'s activations, or must the SAE
itself be re-trained on the adapted model to recover interpretability?
Rather than judging adaptation purely by downstream accuracy, we treat the
encoder's intermediate activations as a representation space and measure
whether the pre-adaptation dictionary continues to describe what the model
computes after adaptation.

\subsection{Problem: what does adaptation do to the SAE basis?}
\label{sec:method-problem}

Let $L$ denote the frozen visual encoder of a VLM, mapping an image
$x \in \mathbb{R}^{H \times W \times 3}$ to patch-token representations
$H^{\ell}(x) \in \mathbb{R}^{N \times d}$ at hook layer $\ell$. Let $A$
denote an adapter (LoRA, MaPLe, or TAP) that produces the adapted encoder
$\hat{L}$. We use $S_L$ to denote an SAE trained on activations from $L$
and $S_{\hat{L}}$ an SAE trained on activations from $\hat{L}$. For an
image $I$, the adapted model produces $\hat{L}(I)$; we then have two
possible sparse decompositions: $S_L(\hat{L}(I))$, which reuses the base
dictionary, and $S_{\hat{L}}(\hat{L}(I))$, which uses the adapted dictionary.

The central question is whether $S_L$ and $S_{\hat{L}}$ agree. If adaptation
merely reweights existing base features (as prompt-tuning work has argued
for near-domain benchmarks), then $S_L \approx S_{\hat{L}}$ and base-SAE
explanations transfer. If adaptation induces new representational structure,
then $S_L(\hat{L}(I))$ will degrade---in reconstruction fidelity, sparsity,
and concept purity---and interpretability claims drawn from $S_L$ will not
carry over to $\hat{L}$. Our method instantiates this comparison and
quantifies the gap.

\begin{itemize}
    \item \textbf{SAEs have become the de facto tool for mechanistic interpretability of LLMs and vision foundation models} \cite{bricken2023monosemanticity, cunningham2024sparse, templeton2024scaling, gao2024scaling, lieberum2024gemma}.

    \item \textbf{All existing pipelines follow the same recipe:} train the SAE once on a large, generic, pretraining-like corpus, freeze it, and reuse it for every downstream task. Representative scales: 8B tokens \cite{bricken2023monosemanticity}, 40B tokens on GPT-4 \cite{gao2024scaling}, and the full Gemma-2 pretraining mix \cite{lieberum2024gemma}. This pattern persists even when the underlying model is adapted via LoRA or other parameter-efficient finetuning: the frozen base-model SAE is typically reused without retraining \cite{kissane2024saes, kharlapenko2024transfer}.

    \item \textbf{This implicitly assumes that deployment activations stay close to training activations}, an assumption inherited from the train-once-freeze-forever pattern and rarely stated.

    \item \textbf{This assumption fails in common deployment regimes.} Activation distributions shift across base $\rightarrow$ instruct/chat finetuning \cite{kissane2024saes, li2025fast}, base $\rightarrow$ domain-specialized finetunes \cite{kharlapenko2024transfer}, pretraining $\rightarrow$ narrow downstream tasks \cite{smith2025negative}, and compositional or OOD concept co-occurrence \cite{menon2026stopprobing}.

    \item \textbf{Empirically, the recipe breaks on exactly these shifts.} On OOD harmful-intent probing, SAE features underperform linear probes on the residual stream, and finetuning the SAE on chat data closes only about half the gap \cite{smith2025negative}. Across data scarcity, class imbalance, label noise, and covariate shift, SAEs show no consistent advantage over simple linear baselines \cite{kantamneni2025sparse}. Under OOD compositional shift, the amortization gap widens and the failure localizes to the learned dictionary itself rather than the inference procedure \cite{menon2026stopprobing}. Independent work frames these findings as evidence that current SAEs produce brittle representations under distribution shift \cite{narayanaswamy2026masked}.

    \item \textbf{Diagnosis:} the failure is structural. A frozen SAE encodes training-time co-occurrence statistics as fixed directions; once the activation distribution shifts, those directions are neither necessary nor sufficient for the deployment domain. A one-shot, pretrain-and-freeze SAE is a model of the training corpus, not of the network.
\end{itemize}

\paragraph{From observations to hypotheses.}
If the dictionary, not the encoder, is the binding constraint, then four
testable claims follow. The first asks whether matching the dictionary
to the target distribution recovers \emph{downstream} signal (H1).
The second asks whether it does so with a qualitatively better
\emph{reconstruction-sparsity} structure (H2).  The third asks whether
the standard \emph{monosemanticity metrics} we would use to adjudicate
H1 and H2 are themselves trustworthy under shift, or whether they need
to be replaced with coverage-aware alternatives (H3).  The fourth asks
whether the base-vs-domain gap \emph{scales with shift severity},
giving a practical deployment criterion (H4).  We state each precisely
below.

\begin{itemize}
    \item \textbf{H1 (Distribution dependence of downstream utility).}
    Sparse autoencoders are distribution-dependent: dictionaries trained on
    the target distribution preserve downstream task performance better
    than those trained on the source distribution. Concretely, a domain
    SAE's encode-decode cycle retains more classification-relevant signal
    than a base SAE's, and zero-shot accuracy drops less under domain-SAE
    intervention across target domains.

    \item \textbf{H2 (Distribution dependence of reconstruction-sparsity structure).}
    Distribution-matched SAEs achieve a more efficient
    reconstruction-sparsity trade-off, encoding target-domain structure
    with fewer but more informative features. Reconstruction loss, $L_0$
    sparsity, and activation entropy jointly improve when the dictionary
    is matched to the target activation space, indicating dictionary
    elements aligned with true modes of variation rather than irrelevant
    source-distribution directions.

    \item \textbf{H3 (Coverage-aware interpretability).}
    Standard monosemanticity metrics overestimate interpretability under
    distribution shift; coverage-aware evaluation reveals that
    distribution-matched SAEs are more monosemantic on the target domain.
    A base SAE can appear spuriously monosemantic by consistently mapping
    unfamiliar inputs to a small set of generic features. Our proposed
    DAMS metric decomposes monosemanticity into Effective Coverage (EC),
    Concept Separability (CSS), and Feature Specificity (FSS), and
    correctly ranks domain SAEs higher across target datasets.

    \item \textbf{H4 (Shift severity governs the gap).}
    The interpretability gap between source-trained and target-trained
    SAEs scales with the magnitude of domain shift, revealing a failure
    boundary for unadapted dictionaries. As pretraining-to-target
    distance increases, base SAEs degrade faster than domain SAEs on all
    metrics. The gap widens monotonically with shift severity, showing
    that dictionaries are distribution-bound objects that require
    updating in proportion to the shift.
\end{itemize}

